
\documentclass[11pt]{article}

\usepackage[a4paper,margin=2.5cm]{geometry}
\usepackage[utf8]{inputenc}
\usepackage[T1]{fontenc}
\usepackage{amsmath,amssymb,amsfonts}
\usepackage{graphicx}
\usepackage{booktabs}
\usepackage{hyperref}
\usepackage{authblk}
\usepackage{bm}
\usepackage{microtype}
\usepackage{xcolor}

\hypersetup{
  colorlinks=true,
  linkcolor=blue!60!black,
  citecolor=blue!60!black,
  urlcolor=blue!60!black
}

\newcommand{\DC}{D_C}
\newcommand{\eps}{\varepsilon}
\newcommand{\Geff}{G_{\rm eff}}
\newcommand{\sig}{\sigma_8}
\newcommand{\LCDM}{\Lambda\mathrm{CDM}}
\newcommand{\BuP}{\textsc{BuP}}
\newcommand{\camb}{\textsc{camb}}

% ─────────────────────────────────────────────────────────────
\title{
  \textbf{Résolution géométrique de la tension~$\sigma_8$\\
  par une dimension cosmologique émergente}
}

\author{Farid Hamdad}

\affil{
  Cadre Bottom-Up Quantum Gravity (\BuP)\\
  \small\url{https://github.com/Farid-Hamdad/Bottom-up-Quantum-Gravity}
}

\date{2026}
% ─────────────────────────────────────────────────────────────

\begin{document}

\maketitle

% ─────────────────────────────────────────────────────────────
\begin{abstract}
Nous montrons que la tension~$\sig$ peut être naturellement atténuée
par une dimension spatiale effective dépendant du temps, dans le cadre
de la gravité quantique bottom-up (\BuP)~\cite{hamdad2026bup}.
La dimension cosmologique évolue selon un profil tardif~$d(z)$
contraint par les données BAO et supernovæ~\cite{desi2024,pantheon},
tout en préservant la physique CMB standard via la dimension critique
asymptotique $\DC = 3.0598 > 3$.
Le couplage gravitationnel effectif
$\Geff(z) = \frac{2}{d(z)-1}\,G$
doit affecter de façon cohérente à la fois l'expansion du fond
et la croissance linéaire des perturbations.
La légère sur-dimensionnalité primordiale ($\DC > 3$) produit une
faible suppression de la croissance des structures à haut redshift
($\Geff < G$), tandis que le regroupement à bas redshift est renforcé
lorsque $d(z) < 3$.
L'intégration complète de l'histoire de croissance depuis la
recombinaison jusqu'à aujourd'hui donne
$\sig = 0.772$,
en accord avec les contraintes de lentillage faible de KiDS-1000~\cite{kids2021}
et DES~Y3~\cite{des2022}, sans introduire de paramètres libres
supplémentaires.
Cela résout l'excès apparent de~$\sig$ obtenu lorsque seule l'évolution
du fond est modifiée, et fournit un mécanisme purement géométrique reliant
la géométrie quantique émergente à la formation des grandes structures.
\end{abstract}

% ─────────────────────────────────────────────────────────────
\section{Introduction}
% ─────────────────────────────────────────────────────────────

L'une des principales tensions de la cosmologie moderne concerne
l'amplitude des fluctuations de matière, quantifiée par le
paramètre~$\sig$.
Les observations CMB Planck dans le cadre~$\LCDM$ standard
favorisent $\sig \simeq 0.811$~\cite{planck2018}, tandis que les
relevés de lentillage faible comme KiDS-1000
($0.766 \pm 0.020$)~\cite{kids2021} et DES~Y3
($0.776 \pm 0.017$)~\cite{des2022} préfèrent des valeurs
systématiquement plus basses.
Cette discordance de $2$--$3\sigma$ constitue la tension~$\sig$.

Dans le cadre \BuP, la géométrie de l'espace-temps n'est pas
fondamentale mais émerge de la structure globale de l'intrication
quantique~\cite{hamdad2026bup}. La dimension spatiale effective est
donc dynamique et peut dévier de la valeur standard $d = 3$.

Des travaux précédents~\cite{hamdad2026bup} ont montré que :
\begin{itemize}
  \item une dimension constante $d \neq 3$ est doublement exclue par
        les contraintes CMB et BAO ($r_{\rm drag}$ chute de~15\%,
        $\sig$ augmente de~62\% pour $d = 2.65$) ;
  \item une dimension variable tardive~$d(z)$ améliore l'ajustement
        aux données DESI~2024 et Pantheon+
        ($\Delta\mathrm{AIC} = -6.23$, $\Delta\mathrm{BIC} = -3.71$) ;
  \item la valeur présente $d(z=0) \sim 2.4$ est cohérente avec les
        courbes de rotation galactiques SPARC.
\end{itemize}

Cependant, modifier seulement l'expansion du fond en conservant la
croissance des perturbations standard conduit à $\sig \sim 0.86$,
aggravant la tension observationnelle.
Dans ce papier, nous montrons que c'est un artefact d'un traitement
incohérent. Une fois que~$\Geff(z)$ est appliqué de façon cohérente
à la croissance des perturbations, la tension disparaît.

% ─────────────────────────────────────────────────────────────
\section{Dimension cosmologique émergente}
% ─────────────────────────────────────────────────────────────

La dimension cosmologique effective est modélisée
comme~\cite{hamdad2026bup}
\begin{equation}
  d(z) = \DC - \frac{\Delta d}{1 + X(z)^\alpha},
  \qquad
  X(z) = X_0\,(1+z)^\beta,
  \label{eq:dz}
\end{equation}
où $X(z)$ est une variable effective de granulation grossière encodant
l'approche vers la dimension critique.

L'exposant $\alpha = 1.78$ n'est pas un paramètre cosmologique libre :
il est fixé de façon indépendante par l'analyse de taille finie (FSS)
de la dimension spectrale sur les graphes d'intrication BuP pour
$N = 6$--$16$ qubits~\cite{hamdad2026bup}.
Les paramètres cosmologiques meilleur ajustement sur DESI~2024 + Pantheon+ sont
\begin{equation}
  X_0 = 0.537, \quad \beta = 2.000, \quad \Delta d = 0.878,
  \quad H_0 = 63.67\;\mathrm{km/s/Mpc}, \quad \Omega_m = 0.3448.
\end{equation}
Cela donne $d(z=0) \simeq 2.40$ et
$d(z\to\infty) \to \DC = 3.0598$,
préservant la physique de recombinaison standard.

% ─────────────────────────────────────────────────────────────
\section{Couplage gravitationnel effectif}
% ─────────────────────────────────────────────────────────────

Le cadre \BuP\ prédit un couplage gravitationnel modifié
géométriquement par la dimension émergente :
\begin{equation}
  \Geff(z) = \frac{2}{d(z)-1}\,G.
  \label{eq:geff}
\end{equation}
Cette relation doit s'appliquer de façon cohérente à la fois à
l'expansion du fond et à la croissance linéaire des perturbations.

\subsection*{Régime à haut redshift}

Aux temps primordiaux, $d(z) \to \DC > 3$, donc $\Geff < G$.
Pour le modèle meilleur ajustement, $\Geff(z \gg 1) \simeq 0.971\,G$, produisant
une légère suppression persistante de la croissance des structures
sur la majeure partie de l'histoire cosmique
(${\sim}99\%$ du temps depuis la recombinaison).

\subsection*{Régime à bas redshift}

Aux temps tardifs, $d(z) < 3$, donc $\Geff > G$.
Aujourd'hui, $d(0) \simeq 2.40$ donne $\Geff(0) \simeq 1.43\,G$,
renforçant le regroupement de matière tardif.
Ce renforcement n'agit que sur ${\sim}1\%$ de l'histoire de
croissance totale.

La compétition entre ces deux régimes est illustrée dans la
Figure~\ref{fig:geff}.

% ─────────────────────────────────────────────────────────────
\section{Croissance linéaire et $\sigma_8$}
% ─────────────────────────────────────────────────────────────

Nous résolvons l'équation de croissance linéaire
\begin{equation}
  D''(a)
  + \left(\frac{3}{a} + \frac{d\ln E}{da}\right) D'(a)
  = \frac{3\Omega_m}{2\,a^5\,E(a)^2}\,\frac{\Geff(z)}{G}\,D(a),
  \label{eq:growth}
\end{equation}
de $a = 10^{-3}$ ($z \sim 1000$) jusqu'à aujourd'hui, avec
$E(z) = H(z)/H_0$ et $\Geff/G$ donné par l'Éq.~(\ref{eq:geff}).
Conditions initiales : $D = a$, $D' = 1$ en $a = 10^{-3}$ (régime
matière-dominé). Normalisation : $\sig = 0.8114 \times
D_{\rm BuP}/D_{\LCDM}$.

La valeur obtenue est
\begin{equation}
  \boxed{\sig = 0.7716}
\end{equation}
à comparer avec
$\sig^{\rm KiDS} = 0.766 \pm 0.020$~\cite{kids2021}
et
$\sig^{\rm DES} = 0.776 \pm 0.017$~\cite{des2022}.
La prédiction BuP est pleinement cohérente avec les deux relevés.

La robustesse de ce résultat est démontrée dans la
Figure~\ref{fig:geff} (panneau droit), qui montre~$\sig$ en fonction
de $d(z=0)$ sur toute la gamme physique $\Delta d \in [0.3, 0.9]$.
Toutes les valeurs restent dans~$1\sigma$ de KiDS-1000.

% ─────────────────────────────────────────────────────────────
\section{Résultats}
% ─────────────────────────────────────────────────────────────

\begin{table}[htbp]
\centering
\caption{Comparaison des prédictions de croissance des structures pour
différents traitements du modèle \BuP. La tension KiDS est calculée
comme $(\sig - 0.766)/0.020$.}
\label{tab:sigma8}
\smallskip
\begin{tabular}{lccc}
\toprule
Modèle & $\sig$ & Tension (KiDS) & Résultat \\
\midrule
$\LCDM$ (Planck~2018) & 0.811 & $+2.3\sigma$ & Tension \\
\BuP\ --- fond seul & 0.861 & $+4.7\sigma$ & Artefact \\
\BuP\ --- $\Geff$ cohérent (ce travail) & \textbf{0.772}
  & $+0.3\sigma$ & Cohérent \\
\bottomrule
\end{tabular}
\end{table}

La valeur précédente $\sig \sim 0.86$ n'était pas une prédiction
authentique du modèle, mais la conséquence d'une approximation
incohérente où seul le fond était modifié.
Une fois $\Geff(z)$ inclus dans la croissance des perturbations,
la tension disparaît.

Cela établit que la légère sur-dimensionnalité primordiale $\DC > 3$
n'est pas une correction mineure, mais le mécanisme physique central
du modèle.

\begin{figure}[htbp]
  \centering
  \includegraphics[width=\textwidth]{fig3_sigma8_geff}
  \caption{\textit{Gauche} : profil de $\Geff(z)/G$ pour le modèle
    \BuP\ (meilleur ajustement). La région bleue indique l'inhibition à haut~$z$
    ($\Geff < G$, dominante sur ${\sim}99\%$ de l'histoire cosmique) ;
    la région rouge le renforcement à bas~$z$ ($\Geff > G$).
    \textit{Droite} : $\sig$ en fonction de $d(z=0)$, avec traitement
    cohérent de~$\Geff$. Les bandes colorées montrent les contraintes
    KiDS-1000 ($\pm 2\sigma$) et DES~Y3 ($\pm 2\sigma$).
    Toutes les valeurs dans la gamme physique
    $\Delta d \in [0.3, 0.9]$ sont cohérentes avec les deux relevés.}
  \label{fig:geff}
\end{figure}

% ─────────────────────────────────────────────────────────────
\section{Discussion}
% ─────────────────────────────────────────────────────────────

Le résultat le plus important est que le même mécanisme géométrique
satisfait simultanément :
\begin{itemize}
  \item la préservation du CMB ($d \to \DC$ à haut~$z$,
        $\theta_s$ inchangé) ;
  \item l'amélioration des ajustements BAO + supernovæ
        ($\Delta\mathrm{AIC} = -6.23$ par rapport à~$\LCDM$) ;
  \item la cohérence avec les courbes de rotation galactiques SPARC
        ($d(z=0) \approx 2.40$, dans~$0.8\sigma$ de SPARC) ;
  \item la résolution de la tension~$\sig$
        ($+0.3\sigma$ par rapport à KiDS-1000).
\end{itemize}
Très peu de modèles de gravité modifiée satisfont les quatre conditions
simultanément sans large liberté paramétrique supplémentaire.

\paragraph{Tension $H_0$ et prédiction JWST.}
La valeur best-fit $H_0 \simeq 63.7$~km/s/Mpc reste inférieure à
Planck ($67.4 \pm 0.5$), une tension partagée avec tous les modèles
d'énergie sombre dynamique contraints sur BAO+SNe.
Elle n'est pas spécifique à \BuP~\cite{tensions2025}.

Une conséquence directe de cette valeur de~$H_0$ est que \BuP\ prédit
un Univers légèrement plus vieux :
\begin{equation}
  t_0 = 14.24\;\text{Gyr}
  \quad
  \text{(contre } 13.80\;\text{Gyr en } \LCDM\text{)},
\end{equation}
soit environ $430$~Myr supplémentaires disponibles pour la formation
des structures à $z > 10$.
Combiné au couplage gravitationnel renforcé $\Geff > G$ aux temps tardifs,
Combiné au couplage gravitationnel renforcé $\Geff > G$ aux temps tardifs,
\BuP\ suggère une contribution possible à l'excès de galaxies massives
observé par JWST~\cite{jwst2023}. Toutefois, une description quantitative
nécessite l'introduction d'inhomogénéités locales $d(z,\delta)$,
étudiées dans un travail séparé~\cite{hamdad2026p4}.



La résolution simultanée des tensions~$H_0$ et~$\sig$ nécessite
probablement soit une modification de la physique de recombinaison,
soit un couplage dépendant de l'échelle $\Geff(z, k)$.


\paragraph{Évaluation quantitative de la prédiction JWST.}
Une analyse numérique dédiée précise la prédiction qualitative ci-dessus.
À $z > 7$, la différence d'âge entre \BuP\ et $\LCDM$ n'est que
$\Delta t \simeq +6$--$9$~Myr, insuffisant pour accélérer significativement
la formation des galaxies via le mécanisme de fond seul. De plus,
$\Geff \simeq 0.971\,G < G$ pour $z > 2$, produisant une légère
suppression de la croissance moyenne des structures. Le $d(z)$ de fond
seul ne prédit donc pas un fort excès de galaxies massives à $z > 7$.

Ce résultat motive une extension naturelle~: une dimension localement
réduite dans les surdensités, $d(z,\delta) = d_{\rm bg}(z) - \eps\,\delta$,
qui accélère l'effondrement dans les pics proto-galactiques rares tout
en préservant le $\sig = 0.772$ global. Ce mécanisme est développé dans
le Paper~4 compagnon~\cite{hamdad2026p4}, où une mesure microscopique
directe sur des graphes d'intrication \BuP\ à $N=9$--$20$ qubits donne
$\eps_{\rm micro}(N=20) = 0.0065$ (IC~68\% $[0.0024, 0.010]$),
remarquablement proche de la valeur calibrée JWST $\eps_{\rm JWST} = 0.011$.

\paragraph{Limites.}
Le calcul actuel utilise l'ODE de croissance linéaire exacte plutôt
qu'une implémentation Boltzmann complète.
Une implémentation Boltzmann complète via \camb/\textsc{class}
incluant l'ensemble du secteur des perturbations constitue la
prochaine étape naturelle de ce programme.
Le résultat $\sig = 0.772$ doit donc être compris comme une estimation
robuste en attente de cette confirmation.

% ─────────────────────────────────────────────────────────────
\section{Conclusion}
% ─────────────────────────────────────────────────────────────

Nous avons montré que la tension~$\sig$ est résolue géométriquement
par une dimension cosmologique émergente.
Le résultat central est $\DC = 3.0598 > 3$, qui implique :
\begin{itemize}
  \item la physique CMB standard aux temps primordiaux
        ($\Geff \to G$ à haut~$z$) ;
  \item une légère suppression de la croissance primordiale des
        structures ($\Geff \simeq 0.971\,G$ sur ${\sim}99\%$ de
        l'histoire cosmique) ;
  \item un renforcement du regroupement tardif
        ($\Geff \simeq 1.43\,G$ aujourd'hui) ;
  \item la prédiction finale $\sig = 0.772$, cohérente avec
        KiDS-1000 et DES~Y3 à~$0.3\sigma$.
\end{itemize}

Une légère sur-dimensionnalité primordiale de l'Univers pourrait
fournir une explication géométrique unifiée pour l'énergie sombre,
la suppression de la croissance des structures, et l'émergence
de la gravité effective elle-même.

% ─────────────────────────────────────────────────────────────
\bibliographystyle{plain}
\begin{thebibliography}{99}

\bibitem{hamdad2026bup}
F.~Hamdad,
\textit{Contraintes cosmologiques sur une dimension spatiale émergente :
dimension variable, énergie noire et $\sigma_8$ dans le cadre
Bottom-Up Quantum Gravity},
Prépublication HAL hal-05590614v1 (2026),
\url{https://hal.science/hal-05590614v1}.

\bibitem{desi2024}
DESI Collaboration,
\textit{DESI 2024 VI : Cosmological Constraints from Baryon Acoustic
Oscillations},
arXiv:2404.03002 (2024).

\bibitem{pantheon}
D.~Brout \textit{et al.},
\textit{The Pantheon+ Analysis: Cosmological Constraints},
Astrophys.\ J.\ \textbf{938}, 110 (2022).

\bibitem{kids2021}
KiDS Collaboration (M.~Asgari \textit{et al.}),
\textit{KiDS-1000 Cosmology: Cosmic shear constraints and comparison
between two point statistics},
Astron.\ Astrophys.\ \textbf{645}, A104 (2021).

\bibitem{des2022}
DES Collaboration,
\textit{Dark Energy Survey Year 3 results: Cosmological constraints
from galaxy clustering and weak lensing},
Phys.\ Rev.\ D \textbf{105}, 023520 (2022).

\bibitem{planck2018}
Planck Collaboration,
\textit{Planck 2018 results VI. Cosmological parameters},
Astron.\ Astrophys.\ \textbf{641}, A6 (2020).

\bibitem{tensions2025}
E.~Di~Valentino \textit{et al.},
\textit{Cosmological tensions in 2025},
Nature Astronomy (2026).

\bibitem{jwst2023}
L.~Labbe \textit{et al.},
\textit{A population of red candidate massive galaxies $\sim$600~Myr
after the Big Bang},
Nature \textbf{616}, 266 (2023).

\bibitem{hamdad2026p4}
F.~Hamdad,
\textit{Nonlinear Collapse from Local Dimensional Reduction in
Bottom-Up Quantum Gravity: A Microscopic Origin for the JWST
Massive Galaxy Excess},
HAL preprint (2026).

\end{thebibliography}

\end{document}
